\documentclass[12pt]{article}
\usepackage{geometry}
\geometry{margin=1.4in}
\usepackage{setspace}
\setstretch{1.4}
\usepackage{amsmath}

\title{Economics as Displacement\\
Money Systems Between Fairness and Extraction}
\author{Diego Rincón}
\date{}

\begin{document}

\maketitle

\begin{abstract}
Money claims to be a neutral exchange medium (ground state $S_0$): you trade value for value, fairly, at equilibrium. The actual money system is extraction (actual state $S^*$): central banks create money and lend it to governments at interest; governments pass this cost to citizens as debt; corporations extract surplus from labor; inflation steals from savers; debt enslaves borrowers. The displacement is enormous. Most people feel it: work never pays enough, costs always rise, debt is inescapable. The cost of maintaining this displacement is paid by debtors and the poor. The cost of return — moving to a fair exchange system — is revolutionary. Central banks and creditors will resist with everything they have.
\end{abstract}

\section{The Ground State: Fair Exchange}

In a fair exchange system, value flows both directions:
\begin{itemize}
\item You produce something. It has value $V$.
\item Someone wants it. They trade you something of equal value $V$.
\item The exchange is fair. Neither party loses.
\item Money is neutral: it measures value but does not create it or destroy it.
\end{itemize}

If I bake a loaf of bread (value: 1 hour of skilled labor, ingredient cost), you trade me something worth 1 hour of labor. We both walk away even. No extraction. No surplus.

In this system:
\begin{itemize}
\item Your wage equals the value you create
\item Prices equal production cost + fair margin
\item Interest (if used) compensates the lender for delayed consumption
\item Debt is contractual: you borrow, you repay at agreed terms, then you're free
\end{itemize}

This is $S_0$: equilibrium exchange.

\section{The Actual System: Extraction}

The actual money system is designed for extraction:

\subsection{Debt-Based Money}

Money is created by banks as debt. When you take out a loan for \$100,000, the bank creates \$100,000 and lends it to you. You owe them \$100,000 + interest.

Where does the money for interest come from? It doesn't exist. It must be created by new loans. The system requires perpetual growth: new debt, new money, new growth. If growth stops, defaults cascade and the system collapses.

This is not a stable exchange system. It is a Ponzi scheme.

\subsection{Inflation}

Central banks expand the money supply to stimulate growth. This dilutes the value of existing money. Savers lose. Borrowers (especially government and corporations) gain.

If inflation is 5\% per year, your \$100 in savings is worth \$95 next year. This is a tax on savers.

\subsection{Wage Suppression}

Your wage is set by competitive pressure, not by the value you create. Millions of workers compete for the same job. Employers set wages as low as possible.

If you create \$100/day in value and are paid \$30/day, the difference (\$70/day) is extracted.

\subsection{Debt Enslavement}

A \$300,000 mortgage at 7\% costs \$720,000 over 30 years. You pay \$420,000 in interest alone. Half your labor goes to banks.

\subsection{Inequality Lock-In}

Those with money extract returns. After 30 years, capital has compounded while labor has not.

\section{Who Maintains the Displacement?}

Central banks, governments, and large corporations benefit from extraction. They maintain it through mythology, selection of defenders, normalization, and punishment of critics.

\section{The Return Cost}

To return to fairness requires:
\begin{enumerate}
\item Debt jubilee (cost: trillions to banks)
\item Wage justice (cost: profit margins to corporations)
\item Inflation reversal (cost: real weight to debtors)
\item Wealth tax (cost: billions to the rich)
\item Bank reform (cost: control and profit to central banks)
\item Political reform (cost: power to creditors)
\end{enumerate}

Each step is resisted. Return is revolutionary.

\section{Three Paths Forward}

\subsection{Path 1: Incremental Fairness}

Cap interest rates, enforce living wages, progressive taxation, stabilize money, regulate banks. The rich give up some extraction.

\subsection{Path 2: Structural Reform}

Replace debt-based money with fiat money created and spent by government (not lent). Interest is zero. Extraction from debt is eliminated.

\subsection{Path 3: Alternative Economics}

Build parallel systems: local currencies, cooperatives, mutual aid, skills exchange. These don't replace the mainstream but offer alternatives.

\section{The Mathematics}

40 years at \$50,000/year = \$2,000,000 gross.

Extraction:
\begin{itemize}
\item Taxes and inflation: \$600,000
\item Mortgage interest: \$420,000
\item Student loans: \$50,000
\item Credit card interest: \$50,000
\item Wage suppression: \$280,000
\item Total: \$1,400,000
\end{itemize}

You keep \$600,000. In a fair system, you keep \$2,000,000.

You live 3.3× better if extraction is removed.

\section{The Path}

The system will not return voluntarily. It will be forced by economic collapse, political revolution, or technological disruption.

Until then:
\begin{enumerate}
\item Understand the extraction
\item Minimize your participation (pay off debt)
\item Build alternatives (local currency, cooperative work)
\item Support reform (vote for fair policies)
\end{enumerate}

The system is displacement. Seeing it clearly is the first step to paying less for it.

\end{document}
